\documentclass[a4paper,12pt]{article}
\usepackage[margin=1in]{geometry}
\usepackage{hyperref}
\usepackage{titlesec}
\usepackage{parskip}

\titleformat{\section}{\large\bfseries}{}{0em}{}[\titlerule]
\titlespacing{\section}{0pt}{12pt}{6pt}

\title{\textbf{B201 Computer Science Lab} \\ Portfolio Report}
\author{Sami Albaroudi}
\date{}

\begin{document}

\maketitle

\section{Introduction}
This report documents the design, implementation, and content of my computer science portfolio website. The portfolio was built as part of the B201 Computer Science Lab module at Gisma University of Applied Sciences, and is intended to showcase my technical skills to potential employers.

\begin{itemize}
    \item Portfolio Website: \href{https://ITZSAMI.github.io}{https://ITZSAMI.github.io}
    \item GitHub Repository: \href{https://github.com/ITZSAMI/ITZSAMI.github.io}{https://github.com/ITZSAMI/ITZSAMI.github.io}
\end{itemize}

\section{Repository Structure}
The GitHub repository is organised as follows:

\begin{itemize}
    \item \texttt{index.html} — The main portfolio webpage containing all sections of the site.
    \item \texttt{style.css} — Custom CSS styles that define the visual appearance of the website.
    \item \texttt{Sami\_CV.pdf} — A professional CV generated using LaTeX and linked from the portfolio.
    \item \texttt{cv.tex} — The LaTeX source code used to generate the CV.
    \item \texttt{README.md} — Documentation describing the repository structure and project overview.
\end{itemize}

\section{Website Structure}
The portfolio is a single-page website divided into five sections:

\begin{itemize}
    \item \textbf{Navigation Bar} — A fixed top navbar allowing smooth scrolling to each section.
    \item \textbf{Hero} — An introduction section displaying my name, title, and location, with a call-to-action button.
    \item \textbf{About} — A brief personal introduction including spoken languages.
    \item \textbf{Projects} — Four project cards showcasing my technical work with descriptions, technology tags, and GitHub links.
    \item \textbf{Skills} — A grid displaying my core technical skills with icons.
    \item \textbf{Contact} — Links to my GitHub, LinkedIn, email, and a downloadable CV.
\end{itemize}

\section{Design Decisions}
Several key design decisions were made during the development of this portfolio.

\textbf{Dark Theme:} A dark colour scheme was chosen as it is widely used in developer portfolios and aligns with the aesthetic expectations of the tech industry. It also reduces eye strain for users viewing the site in low-light environments.

\textbf{Bootstrap 5:} Bootstrap was selected as the CSS framework because it provides a responsive grid system and pre-built components that significantly speed up development. This is a deliberate and principled choice, as using established frameworks is standard practice in professional web development.

\textbf{Single Page Layout:} The portfolio uses a single-page design with anchor links for navigation. This was chosen because a portfolio site has limited content that does not warrant multiple pages, and a single page provides a smoother user experience.

\textbf{Bootstrap Icons:} Icons were added using the Bootstrap Icons library to improve visual clarity and make sections easier to scan quickly.

\textbf{GitHub Pages:} The site is hosted on GitHub Pages because it is free, integrates directly with the repository, and automatically deploys on every push to the main branch.

\section{Tools and Technologies}
The following tools and technologies were used to develop the portfolio:

\begin{itemize}
    \item \textbf{HTML5} — Used to structure the content of the webpage.
    \item \textbf{CSS3} — Used for custom styling beyond what Bootstrap provides.
    \item \textbf{Bootstrap 5} — Used for responsive layout and pre-built UI components.
    \item \textbf{Bootstrap Icons} — Used for iconography throughout the site.
    \item \textbf{Git \& GitHub} — Used for version control and hosting the repository.
    \item \textbf{GitHub Pages} — Used to deploy and host the live website.
    \item \textbf{LaTeX (Overleaf)} — Used to create the CV and this report.
    \item \textbf{VS Code} — Used as the primary code editor.
\end{itemize}

\section{Strengths and Weaknesses}

\textbf{Strengths:}
\begin{itemize}
    \item The site is fully responsive and works on both desktop and mobile devices.
    \item The design is clean, modern, and appropriate for a developer portfolio.
    \item The repository is well structured and documented with a clear README.
    \item The CV is professionally typeset using LaTeX and available for download directly from the site.
\end{itemize}

\textbf{Weaknesses:}
\begin{itemize}
    \item The portfolio currently has no backend or dynamic functionality.
    \item Project descriptions are brief and could be expanded with more technical detail.
    \item There is no dedicated project detail page for each project.
\end{itemize}

\section{Future Improvements}
In the future, I would like to make the following improvements to the portfolio:

\begin{itemize}
    \item Add a contact form with a backend so visitors can send messages directly.
    \item Create individual project pages with screenshots, technical explanations, and live demos.
    \item Add animations and scroll effects to improve the user experience.
    \item Expand the projects section as I build more technical projects during my studies.
    \item Rebuild the site using React to demonstrate more advanced frontend skills relevant to Werkstudent roles in Berlin.
\end{itemize}

\section{Conclusion}
This portfolio represents my current technical skills and serves as a practical demonstration of my ability to design, implement, and deploy a web-based project. The development process involved making principled decisions about technology choices, visual design, and content structure. I plan to continue improving the portfolio as I develop new skills throughout my studies.
\section{References}

\begin{itemize}
    \item Bootstrap (2023) \textit{Bootstrap 5 Documentation}. Available at: \href{https://getbootstrap.com}{https://getbootstrap.com} (Accessed: 26 June 2026).
    
    \item Bootstrap Icons (2023) \textit{Bootstrap Icons Documentation}. Available at: \href{https://icons.getbootstrap.com}{https://icons.getbootstrap.com} (Accessed: 26 June 2026).
    
    \item GitHub (2023) \textit{GitHub Pages Documentation}. Available at: \href{https://pages.github.com}{https://pages.github.com} (Accessed: 26 June 2026).
    
    \item Overleaf (2023) \textit{Overleaf Documentation}. Available at: \href{https://www.overleaf.com/learn}{https://www.overleaf.com/learn} (Accessed: 26 June 2026).
\end{itemize}
\end{document}